\documentclass[12pt,a4paper]{memoir}
% \documentclass[titlepage,12pt,a4paper]{book}

% substituir linha seguinte por 
% \usepackage[english]{babel} 
% se o relatório for escrito na língua inglesa.
\usepackage[portuguese]{babel}

% \usepackage[utf8]{inputenc}
\usepackage[T1]{fontenc}

\usepackage{makeidx}
\usepackage{xspace}
\usepackage{graphicx,color,times}
\usepackage{fancyhdr}
% \usepackage{pxfonts}
% \usepackage{times}
% \usepackage{mathptm}
% \usepackage{amssymb}
% \usepackage{amsfonts}

\usepackage{amsmath}
\usepackage{latexsym}
\usepackage[printonlyused]{acronym}
\usepackage{float}
\usepackage{listings}
\usepackage{tocbibind}
\usepackage{natbib}
\usepackage{hyperref}

% \usepackage{glossaries}
% \makeglossaries

% \renewcommand{\ttdefault}{phv}

\pagestyle{fancy}
\renewcommand{\chaptermark}[1]{\markboth{#1}{}}
\renewcommand{\sectionmark}[1]{\markright{\thesection\ #1}}
\fancyhf{} \fancyhead[LE,RO]{\bfseries\thepage}
\fancyhead[LO]{\bfseries\rightmark}
\fancyhead[RE]{\bfseries\leftmark}
\renewcommand{\headrulewidth}{0.5pt}
\renewcommand{\footrulewidth}{0pt}
\setlength{\headheight}{15.6pt}
\setlength{\marginparsep}{0cm}
\setlength{\marginparwidth}{0cm}
\setlength{\marginparpush}{0cm}
\addtolength{\hoffset}{-1.0cm}
\addtolength{\oddsidemargin}{\evensidemargin}
\addtolength{\oddsidemargin}{0.5cm}
\addtolength{\evensidemargin}{-0.5cm}


\usepackage{fix-cm}
\usepackage{fourier}
\usepackage[scaled=.92]{helvet}
\definecolor{ChapGrey}{rgb}{0.6,0.6,0.6}
\newcommand{\LargeFont}{
  \usefont{\encodingdefault}{\rmdefault}{b}{n}
  \fontsize{60}{80}\selectfont\color{ChapGrey}
  }
\makeatletter
\makechapterstyle{GreyNum}{
  \renewcommand{\chapnamefont}{\large\sffamily\bfseries\itshape}
  \renewcommand{\chapnumfont}{\LargeFont}
  \renewcommand{\chaptitlefont}{\Huge\sffamily\bfseries\itshape}
  \setlength{\beforechapskip}{0pt}
  \setlength{\midchapskip}{40pt}
  \setlength{\afterchapskip}{60pt}
  \renewcommand\chapterheadstart{\vspace*{\beforechapskip}}
  \renewcommand\printchaptername{
  \begin{tabular}{@{}c@{}}
    \chapnamefont \@chapapp\\}
    \renewcommand\chapternamenum{\noalign{\vskip 2ex}}
    \renewcommand\printchapternum{\chapnumfont\thechapter\par}
    \renewcommand\afterchapternum{
  \end{tabular}
  \par\nobreak\vskip\midchapskip}
  \renewcommand\printchapternonum{}
  \renewcommand\printchaptertitle[1]{
  {\chaptitlefont{##1}\par}}
  \renewcommand\afterchaptertitle{\par\nobreak\vskip \afterchapskip}
}
\makeatother
\chapterstyle{GreyNum}

\setcounter{tocdepth}{3}
\setsecnumdepth{subsubsection}

\renewcommand{\ttdefault}{lmtt}


% NEW COLORS
\definecolor{dark}{gray}{0.25}
\definecolor{lgray}{gray}{0.9}
\definecolor{dkblue}{rgb}{0,0.13,0.4}
\definecolor{dkgreen}{rgb}{0,0.6,0}
\definecolor{gray}{rgb}{0.5,0.5,0.5}
\definecolor{mauve}{rgb}{0.58,0,0.82}

\lstset{ %
  language=C,                    basicstyle=\footnotesize,
  numbers=none,                  numberstyle=\tiny\color{gray}, 
  stepnumber=1,                  numbersep=5pt,
  backgroundcolor=\color{white}, showspaces=false,
  showstringspaces=false,        showtabs=false,
  frame=single,                  rulecolor=\color{black},
  tabsize=2,                     captionpos=b,
  breaklines=true,               breakatwhitespace=false,
  title=\lstname,                keywordstyle=\color{blue},
  commentstyle=\color{dkgreen},  stringstyle=\color{mauve},
  escapeinside={\%*}{*)},        morekeywords={*},
  belowskip=0cm
}

\renewcommand{\lstlistingname}{Excerto de Código}
\renewcommand{\lstlistlistingname}{Lista de Excertos de Código}

\renewcommand{\today}{\day \ifcase \month \or Janeiro\or Fevereiro\or Março\or %
Abril\or Maio\or Junho\or Julho\or Agosto\or Setembro\or Outubro\or Novembro\or %
Dezembro\fi de \number \year} 



\begin{document}


\thispagestyle{empty}
\setcounter{page}{-1}

\begin{center}
\begin{Huge}
\textbf{Universidade da Beira Interior}
\end{Huge}
\end{center}

\begin{center}
\begin{Huge}
Departamento de Informática
\end{Huge}
\end{center}

\vspace{0,07cm}
\begin{figure}[!htb]
\centering
\includegraphics[width=191pt]{ubi-fe-di.png}
\end{figure}

\vspace{0.5cm}
\begin{center}
\begin{Large}
\textbf{2020/2021: \emph{SISTEMA SOLAR}}
\end{Large}
\end{center}


\vspace{0.5cm}
\begin{center}
\begin{normalsize}
\begin{large}
Elaborado por:
\end{large}
\end{normalsize}
\end{center}


\begin{center}
\begin{large}
\textbf{Cristiano Santos, nº 43464}


\textbf{Bruno Monteiro, nº 43994}


\textbf{Alexandre Monteiro, nº 44149}
\end{large}
\end{center}

\vspace{0.2cm}
\begin{center}
\begin{large}
\textbf{Unidade Curricular: Computação Gráfica}
\end{large}
\end{center}

\vspace{0,3cm}
\begin{center}
\begin{normalsize}
\begin{large}
Docentes da Cadeira:
\end{large}
\end{normalsize}
\end{center}

\begin{center}
\begin{large}
\textbf{Professor Doutor Abel João Padrão Gomes}
\end{large}
\end{center}


\begin{center}
\begin{large}
\textbf{Professor Nuno Filipe Alexandre Carapito}
\end{large}
\end{center}

\vspace{0.5cm}
\begin{center}
\begin{normalsize}
\today
\end{normalsize}
\end{center}


\clearpage{\thispagestyle{empty}\cleardoublepage}

\frontmatter

\chapter*{Agradecimentos}
\label{chap:ack}

Queriamos começar por agradecer ao Professor Doutor Abel Gomes por toda a matéria lecionada e pela prosposta de projeto que nos permitiu reforçar os conhecimentos pedagógicos da Unidade Curricular.

Gostávamos também de agradecer ao Professor Nuno Carapito por toda a disponibilidade e ajuda prestada.

\clearpage{\thispagestyle{empty}\cleardoublepage}


\tableofcontents

\clearpage{\thispagestyle{empty}\cleardoublepage}

\listoffigures

% #   ATENÇÃO
% Se não existirem tabelas, comentar as duas linhas seguintes
%\clearpage{\thispagestyle{empty}\cleardoublepage}
%\listoftables

% #   ATENÇÃO
% Se existirem trechos de código, descomentar as seguintes linhas
% \clearpage{\thispagestyle{empty}\cleardoublepage}
% \lstlistoflistings

\clearpage{\thispagestyle{empty}\cleardoublepage}
\chapter*{Acrónimos}
\label{chap:acro}

% #   ATENÇÃO
% A lista de acrónimods deve ser ordenada alfanumericamente.
% Estrangeirismos devem ser realçados em itálico.
% Se o relatório for escrito em Inglês, uma palavra portuguesa é um estrangeirismo.

% O maior acrónimo deve ser colocado neste ponto (reparar que CFIUTE é maior que TCP!).
%               vvvvvv
\begin{acronym}[CFIUTE]

  \acro{OpenGL}{Open Graphics Library}
  \acro{GLEW}{OpenGL Extension Wrangler Library}
  \acro{GLFW}{Graphics Library Framework}
  \acro{API}{Application Programming Interface}
  \acro{GLM}{Generalized Linear Model}
  \acro{IDE}{Integrated Development Environment}
  \acro{UBI}{Universidade da Beira Interior}
\end{acronym}

% \clearpage{\pagestyle{empty}\cleardoublepage}
% \chapter*{Glossário}
\makeglossaries

\newglossaryentry{.NET Framework}
{
  name={.NET Framework},
  description={É uma plataforma para desenvolvimento e funcionamento de aplicações desenvolvida pela Microsoft.}
}

\newglossaryentry{ASP.NET}
{
  name={ASP .Net},
  description={É uma plataforma da Microsoft para o desenvolvimento de aplicações Web e é o sucessor da tecnologia ASP.}
}

\newglossaryentry{CS}
{
  name={C\#},
  description={Lê-se \textit{C Sharp} e é uma linguagem de programação orientada a objectos, desenvolvida pela Microsoft, inicialmente para a plataforma .NET. O C\# é inspirado na junção entre as linguagens C++ e Java.}
}


\newglossaryentry{Java}
{
  name={JAVA},
  description={É uma linguagem de programação orientada a objectos, desenvolvida pela Sun Microsystems na década de 90. Hoje pertence à empresa Oracle.}
}


\newglossaryentry{OpenDMTP}
{
  name={OpenDMTP},
  description={\textit{Open Device Monitoring and Tracking Protocol} é um protocolo e uma \textit{framework} abertos que permite a comunicação bidireccional entre servidores e clientes através da internet.}
}


\newglossaryentry{OpenGTS}
{
  name={Open GTS},
  description={É o primeiro projecto \textit{Open Source} \textit{Web-Based} para controlo de frotas por GPS.}
}


\newglossaryentry{VS2010}
{
  name={Visual Studio 2010},
  description={\textit{Microsoft Visual Studio 2010} é um sistema de desenvolvimento desenvolvido pela Microsoft e é dedicado ao Framework .NET, que contem um conjunto de ferramentas de desenvolvimento projectadas para auxiliar os programadores a enfrentarem desafios complexos.}
}


\newglossaryentry{WebS}
{
	name={Web Service},
	description={Web services são aplicações modulares auto-descritas e auto-contidas, que permitem a integração de sistemas e a comunicação entre aplicações de diferentes tipos.}
}


\newglossaryentry{WebBased}
{
	name={Web Based},
	description={Aplicação desenvolvida para a Web.}
}

\newglossaryentry{Roaming}
{
	name={Roaming},
	description={Define a possibilidade de um utilizador de uma determinada rede obter rede/conecção fora da área geográfica onde foi registado.}
}


\newglossaryentry{Smartphone}
{
	name={Smartphone},
	description={Smartphone é um telefone móvel que contem muitas das principais tecnologias de comunicação e serviços que existem nos computadores pessoais, como acesso a e-mails, serviços de mensagens instantâneas, internet, GPS, entre outros.}
}

\newglossaryentry{TCPIP}
{
	name={TCP/IP},
	description={É um conjunto de protocolos de comunicação entre computadores ligados rede. O nome TCP/IP surge da união entre dois protocolos: o TCP (Transmission Control Protocol) e o protocolo IP (Internet Protocol).}
}

\newglossaryentry{Firewall}
{
	name={Firewall},
	description={É o nome criado para definir um dispositivo para uma rede de computadores que tem como objectivo criar uma política de segurança num determinado ponto de controlo da rede.}
}

\newglossaryentry{JavaScript}
{
	name={JavaScript},
	description={É uma linguagem de programação baseada na linguagem de programação ECMAScript. Actualmente é a linguagem de programação mais utilizada em \textit{``Client-Side''} nos \textit{browsers}.}
}

\newglossaryentry{Flash}
{
	name={Flash},
	description={Desenvolvido pela Macromedia, o Flash é um software utilizado para criação de animações interactivas que funcionam incorporadas em \textit{Browsers}, \textit{Desktop}, \textit{Smartphones}, \textit{Tablets}, e Televisores.}
}


\newglossaryentry{StoredProcedure}
{
	name={Stored Procedure },
	description={É o nome dado a um conjunto de comandos numa base de dados de forma a simplificar a sua utilização.}
}

\newglossaryentry{SQLS}
{
	name={SQL Server 2008},
	description={É um sistema de gestão de base de dados relacional criado pela Microsoft.}
}

\newglossaryentry{Firm}
{
	name={Firmware},
	description={É o conjunto de instruções operacionais programadas directamente no \textit{hardware} de um equipamento electrónico.}
}

\newglossaryentry{browser}
{
	name={Browser},
	description={É um programa de computador que possibilita aos utilizadores uma interacção com documentos virtuais da Internet, também conhecidos como páginas Web.}
}



\clearpage{\thispagestyle{empty}\cleardoublepage}

\mainmatter
\acresetall
\chapter{Introdução}
\label{chap:intro}

\section{Enquadramento}
\label{sec:amb} % CADA SECÇÃO DEVE TER UM LABEL
% CADA FIGURA DEVE TER UM LABEL
% CADA TABELA DEVE TER UM LABEL

Neste trabalho foi nos dada a escolha de um projecto de entre 5 temas disponíveis. O tema escolhido pelo grupo foi o  modelo do Sitema Solar. O trabalho foi realizado em C++ com recusro às bibliotecas \ac{OpenGL}, \ac{GLFW}, \ac{GLEW} e \ac{GLM}.

Este relatório foi feito no contexto da unidade curricular de Computação Gráfica da \ac{UBI}. Foi na \ac{UBI} que desenvolvemos todo o trabalho. 

\subsection{Enquadramento Histórico}
O Sistema Solar foi formado à cerca de 4,5 mil milhões de anos a partir de uma densa nuvem de gás interestelar e poeira.

Quando essa nuvem de poeira entrou em colapso, formou uma nebulosa solar (um disco giratório de material).
Vários aglomerados chocaram e formaram objetos cada vez maiores que se juntaram e deram origem aos planetas, planetas anões e luas.

Oficialmente o Sistema Solar tem 8 planetas e 150 luas conhecidas. Mercúrio, Vénus, Terra e Marte são planetas rochosos, Júpiter e Saturno são gigantes gasosos e Urano e Neptuno são gigantes de gelo.

\section{Motivação}
\label{sec:mot}

A Computação Gráfica (CG) é uma área da Ciência da Computação que se dedica
ao estudo e desenvolvimento de técnicas e algoritmos para a síntese de imagens através do computador. Atualmente, é uma das áreas de maior expansão e importância que propicia o desenvolvimento de trabalhos multidisciplinares.

Assim o projeto que se apresenta ajudará a aprofundar o conhecimento nesta área cada vez mais importante e adicionalmente aprofundar os conhecimentos aprendidos na Unidade Curricular

\section{Objetivos}
\label{sec:obj}

O projecto consiste na realização de um modelo interativo do Sistema Solar, em 2D e 3D, onde  projetamos os astros constituintes do mesmo e os seus movimentos (rotações e translações) a uma escala que consideramos adequada, de modo a transmitir ao utilizador uma ideia correta e concreta da realidade.

Caraterísticas da aplicação Gráfica:
\begin{enumerate}
    \item Modelar o sol, os planetas e os seus respetivos satélites;
    \item Texturizar os planetas par aumentar o realismo;
    \item Menus que permitam obter informação sobre os elementos do Sistema Solar;
    \item A câmera pode mover-se para qualquer posição do Sistema Solar, sendo que este movimento afecta a luz que inside sobre os elementos do Sistema Solar.
    \item Utilização do teclado para fazer zoom ao sistema solar e para alterar a vista.
    \item Iluminação através de um foco de luz proveniente do Sol.
    \item Calcular sombras que os planetas e os satélites poderão originar entre si.
\end{enumerate}    

\section{Organização do Documento} -- 
\label{sec:organ}
% !POR EXEMPLO!
De modo a refletir o trabalho que foi feito, este documento encontra-se estruturado da seguinte forma:
\begin{enumerate}
\item O primeiro capítulo -- \textbf{Introdução} -- apresenta o projeto, a motivação para a sua escolha, o enquadramento para o mesmo, os seus objetivos e a respetiva organização do documento.
\item O segundo capítulo -- \textbf{Tecnologias Utilizadas} -- descreve os conceitos mais importantes no âmbito deste projeto, bem como as tecnologias utilizadas durante do desenvolvimento do projeto.
\item O terceiro capítulo -- \textbf{Etapas de Desenvolvimento} -- Neste capítulo são apresentadas quais as etapas tidas em conta no desenevolvimento do projeto e qual o membro do grupo que as realizou.
\item O quarto capítulo -- \textbf{Descrição do funcionamento do Software} -- Nesta etapa explicamos ao pormenor como funcionam as funcões que integram o código, qual o seu intuito e de que forma se relacionam.
\item O quinto capítulo -- \textbf{Considerações Finais} -- Neste capítulo elaborámos uma conclusão sobre o projeto e o conhecimento retido com a sua elaboração.
\item O sexto capítulo -- \textbf{Referências bibliográficas} -- Lista das fontes e referências bibliográficas utilizadas.
\end{enumerate}

\clearpage{\thispagestyle{empty}\cleardoublepage}
\chapter{Tecnologias Utilizadas}
% OU \chapter{Trabalhos Relacionados}
% OU \chapter{Engenharia de Software}
% OU \chapter{Tecnologias e Ferramentas Utilizadas}
\label{chap:estado-da-arte}

\section{Introdução}
\label{chap2:sec:intro}
Neste capítulo descrevemos  os  conceitos mais  importantes  no âmbito  deste  projeto,  bem  como  as  tecnologias utilizadas durante do desenvolvimento da projeto. 

\section{Conceitos importantes}
\label{chap2:sec:citacoes}

Alguns dos conceitos mais importantes:
\begin{enumerate}
    \item \textbf{Básicos de Geometria} -- Conhecimmentos teóricos sobre vetores, matrizes e algebra;
    \item \textbf{Transformações Geométricas} -- Noções básicas de transformações de matrizes ( matrizes de tranformação, rotação e \textit{scaling});
    \item \textbf{Windows and Viewports} -- Uma viewport define em coordenadas normalizadas uma área retangular no display onde a imagem dos dados aparece. Uma window define a área retangular em coordenadas.
    \item \textbf{Shading and Ilumination} -- O shading é a implementação do modelo de iluminação nos pixeis e superfícies poligonais dos objetos gráficos. Por sua vez, a Ilumination é a forma como o shading é aplicado ns diferentes objetos.
    \item \textbf{Color} -- Atribuição de diferentes cores da escala RGB (em OpenGL, entre 0 e 1) aos diversos pixeis e objetos.
\end{enumerate}

\section{Tecnologias Utilizadas}
\label{chap2:sec:...}
O C++ é uma linguagem de programação e de uso geral. É uma das linguagens mais populares desde 1990 e foi criada por Bjarne Stroustrup.

O OpenGL é uma \ac{API} livre utilizada na computação gráfica para o desenvolvimento de aplicações gráficas, ambientes 3D, jogos, etc.

O GLM generaliza a regressão linear permitindo que este seja relacionado à variável de resposta por meio de uma função de ligação, permitindo que a magnitude da variância de cada medição seja uma função do seu valor previsto.

O GLFW é uma biblioteca que pode ser usada com OpenGL. Esta permite que programadores possam criar e manusear janelas e contextos OpenGL. Permite também interagir com joysticks, rato e teclado.

O GLEW é uma biblioteca de carregamento de extensões C/C++ de código \textit{open-source}. Este fornece mecanismos eficientes para determinar quais as extensões OpenGL que são suportadas na plantaforma alvo.

O Visual Studio \ac{IDE} é uma plataforma de lançamento criativa que se pode usar para alterar, depurar, construir código e publicar um aplicação.

\clearpage{\thispagestyle{empty}\cleardoublepage}
\chapter{Etapas de desenvolvimmento}
% OU \chapter{Trabalhos Relacionados}
% OU \chapter{Engenharia de Software}
% OU \chapter{Tecnologias e Ferramentas Utilizadas}
\label{chap:tecno-ferra}

\section{Introdução}
\label{chap3:sec:intro}
Neste capítulo são apresentadas quais as etapas tidas em conta no desenevolvimento do projeto e qual o membro do  grupo que as realizou.

\section{Parte técnica}
\label{chap3:sec:...}

\begin{enumerate}
    \item \textbf{Modelação} -- Número de objetos que iremos precisar para fazer o sistema solar definir um tamanho para cada planeta e lua. Pensar como iremos realizar a trajetória dos planetas (1);
    \item \textbf{Blender} -- Fazer no Blender planetas e luas (2);
    \item \textbf{Interação} -- Rato responsável pelo movimento da câmera. Teclado(WASD) responsável pelo movimento pela cena (3);
    \item \textbf{Iluminação} --  Pesquisa do melhor tipo de iluminação a usar e sua aplicação (iluminação flat) (4);
    \item \textbf{Texturização} -- Pesquisar e aplicar as texturas dos planetas e luas (5).
\end{enumerate}

\section{Funcionalidades Extras}

\begin{enumerate}
    \item \textbf{Implementação do Menu} -- Constituido pelas opções FreeWalk e "Planetas", sendo esta última, uma opção que permite ao utilizador escolher um planeta e ver um pouco das suas caraterísticas (6);
    \item \textbf{Background} -- Aplicação de um background do universo ao modelo (7).
\end{enumerate}


\section{Relatório}
\label{chap3:sec:concs}

\begin{enumerate}
    \item \textbf{Pesquisa} -- Pesquisa sobre o funcionamento e aplicação do \LaTeX (8);
    \item \textbf{Motivação do trabalho e tecnologias utilizadas} (9);
    \item \textbf{Explicação do Código} (10);
    \item \textbf{Conclusão e considerações finais} -- Breve conclusão sobre o projeto e pequena sintese sobre os conhecimentos consolidados e/ou adquiridos (11);
    \item \textbf{Finalização do relatório} (12).
\end{enumerate}

\section{Atribuição de Tarefas}

\begin{itemize}
    \item \textbf{Cristiano Santos} --  resolução dos tópicos 1, 3, 5, 6, 7, 8, 9, 10, 12;
    \item \textbf{Bruno Monteiro} -- resolução dos tópicos 1, 2, 4, 5, 6, 8, 9, 10, 12;
    \item \textbf{Alexandre Monteiro} -- resolução  dos tópicos 1, 4, 5, 6, 8, 10, 11, 12.
\end{itemize}


\clearpage{\thispagestyle{empty}\cleardoublepage}
\chapter{Descrição do funcionamento do Software}
\label{chap:imp-test}

\section{Introdução}
\label{chap4:sec:intro}
Nesta etapa explicámos ao pormenor como funcionam as funcões que integram o código, qual o seu intuito e de que forma se relacionam.

No desenvolvimento do projeto, foi necessário organizar o código em diferentes ficheiros de modo a torná-lo mais organizado e percetível para quem o interpreta.
Sendo assim, temos os segunites ficheiros: \textit{main.cpp}, \textit{draw.hpp} e o  \textit{transferDataToGPU.hpp}. Por sua vez os \textit{shadders} utilizados foram os seguintes: \textit{TransformVertexShader.vertexshader} e o  \textit{TextureFragmentShader.fragmentshader}.

\section{Funcionamento Geral}
Quando o programa é iniciado, o utilizador através das arrow keys up and down, pode selecionar uma das 3 opções disponíveis: "Free Walk", "Planets" e "Quit". Carregando no enter entra numa das opções.

Na opção "Free Walk" o utilizador pode andar livremente pela cena, usando as teclas "WASD" e o rato. Clicando no "R" tem uma perspetiva real do Sistema Solar e clicando no "F" os planetas ficam mais próximos uns dos outros e mais estagnados para permitir ao utilizador ver os planetas de forma mais pormenorizada. Clicando no "Q" volta ao Menu Inicial.

Na opção "Planets" o utilizador pode usar as arrow keys up and down para trocar entre a história dos diferentes planetas. Tal como no menu anterior, pode usar as teclas "WASD" e o rato para se mover na cena. Clicando no "Q" volta ao Menu Inicial.

Na opção "Quit" o programa é encerrado.

\section{\textit{Main.cpp}}
\label{chap4:sec:...}

Iniciamos este ficheiro com a importação das diferentes bibliotecas necessárias:
\begin{lstlisting}[caption=\textit{Includes} das bibliotecas usadas.]
// Include standard headers
#include <stdio.h>
#include <stdlib.h>
#include <time.h>
#include <vector>

// Include GLEW
#include <GL/glew.h>

// Include GLFW
#include <GLFW/glfw3.h>
GLFWwindow* window;

// Include GLM
#include <glm.hpp>
#include <gtc/matrix_transform.hpp>
using namespace glm;

#include "common/shader.hpp"
#include "common/texture.hpp" 
#include "common/controls.hpp" 
#include "common/objloader.hpp"
#include "common/shader.cpp"
#include "common/texture.cpp" 

#include "common/controls.cpp" 
#include "common/objloader.cpp"

#include "transferDataToGPU.hpp"
#include "draw.hpp"
\end{lstlisting}



Criámos uma \textit{typedef struct INFO}, usada para cada um dos seguintes objetos: Background, Sun, Mercury, Venus, Earth, Mars, Jupiter, Saturn, Uranus, Neptune, Ring, Asteroid, Moon, Menu1, Menu2, Menu3, MercuryHistory, VenusHistory, EarthHistory, MarsHistory, JupiterHistory, SaturnHistory, UranusHistory, NeptuneHistory, AsteroidHistory, MoonHistory, SunHistory.
\begin{lstlisting}[caption= \textit{typedef struct INFO}.]
typedef struct {
	GLuint VertexArrayID;
	GLuint MatrixID;
	GLuint Texture;
	GLuint TextureID;
	GLuint uvbuffer;
	GLuint vertexbuffer;
}INFO;
\end{lstlisting}

Criámos uma \textit{typedef struct VECTORS}, usada para cada um dos seguintes objetos: Sphere, Back, VRing, VAsteroid, Menu;
\begin{lstlisting}[caption= \textit{typedef struct VECTORS}.]
typedef struct {
	std::vector<glm::vec3> vertices;
	std::vector<glm::vec2> uvs;
	std::vector<glm::vec3> normals;
}VECTORS;
\end{lstlisting}

Criámos uma \textit{typedef struct TRANSFORMATIONS}, usada para cada um dos seguintes objetos: transfBackground, transfSun, transfMercury, transfVenus, transfEarth, transfMars, transfJupiter, transfSaturn, transfUranus, transfNeptune, transfRing, transfAsteroid, transfMoon, transfList, transfHistory, randomAsteroid[100];
\begin{lstlisting}[caption= \textit{typedef struct TRANSFORMATIONS}.]
typedef struct {
	glm::mat4 self_rotation;
	glm::mat4 center_rotation;
	glm::mat4 translation;
	glm::mat4 scaling;
	float xPosition;
	float yPosition;
	float zPosition;
}TRANSFORMATIONS;
\end{lstlisting}

Criámos a \textit{typedef struct planetRotation}, que foi necessária para a criação da \textit{typedef struct PLANETSROTATIONS};
\begin{lstlisting}[caption= \textit{typedef struct planetRotation}.]
typedef struct {
	float center_Rotation;
	float self_Rotation;
}planetRotation;
\end{lstlisting}

Criámos uma \textit{typedef struct PLANETSROTATIONS}, usada para cada um dos seguintes objetos: UsingRotation;
\begin{lstlisting}[caption= \textit{typedef struct PLANETSROTATIONS}.]
typedef struct {
	planetRotation sun;
	planetRotation mercury;
	planetRotation venus;
	planetRotation earth;
	planetRotation mars;
	planetRotation jupiter;
	planetRotation saturn;
	planetRotation uranus;
	planetRotation neptune;
	planetRotation moon;
	planetRotation asteroidBelt;
}PLANETSROTATIONS;
\end{lstlisting}

Criámos uma \textit{typedef struct DIMENSIONS}, usada para cada um dos seguintes objetos: Real, Fake, Using;
\begin{lstlisting}[caption= \textit{typedef struct DIMENSIONS}.]
typedef struct {
	float sun;
	float mercury;
	float venus;
	float earth;
	float mars;
	float jupiter;
	float saturn;
	float uranus;
	float neptune;
	float moon;
	float asteroidBelt;
}DIMENSIONS;
\end{lstlisting}

Após a criação das \textit{typedef structs}, referenciamos as funções criadas e utilizadas.
São elas as seguintes:
\begin{itemize}
\item \textit{transferDataToGPU()}, contida no \textit{transferDataToGPU.hpp};
\item \textit{draw(), setMVP()}, contida no \textit{draw.hpp};
\item \textit{resetPosPlanetsList()};
\item \textit{selectingMenu1()};
\item \textit{selectingMenu3()};
\item \textit{declareDimensions()};
\item \textit{incrementRotations()};
\item \textit{CelestialBodiesPosition()};
\item \textit{deletebuffers()};
\end{itemize}

\subsection{Funções Implementadas}
\subsubsection{\textit{transferDataToGPU()}}
Começámos por passar como parâmetros a struct INFO, Vectors, a variavel associada aos shaders, o nome da textura e do objeto.

\begin{lstlisting}[caption=\textit{VertexArray} funções.]
glGenVertexArrays(1, (&a->VertexArrayID));
glBindVertexArray((a->VertexArrayID));
\end{lstlisting}
Nestas duas funções vamos buscar o número de Arrays que existem e atribuimo-los ao VertexArrayID da nossa struct INFO.

\begin{lstlisting}[caption=\textit{VertexArray} funções.]
	a->MatrixID = glGetUniformLocation(programID, "MVP");

	a->Texture = loadDDS(textureName);

	a->TextureID = glGetUniformLocation(programID, "myTextureSampler");

	bool res = loadOBJ(objectName, Sphere->vertices, Sphere->uvs, Sphere->normals);
\end{lstlisting}

A primeira função vai utilizar o \textit{MatrixID} para lhe atribuir a localizaçao da matriz MVP dos \textit{shaders}. A segunda vai atribuir a variável textura à nossa textura do que nos é passada em parametro para uma função de carregar texturas (loadDDS(\textit{texturename()}). A terceira textura, no fundo, faz o mesmo que a priemira função mas aplicada à textura. A quarta e ultima, vai atribuir à nossa \textit{struct VECTORS} os vértices, os UVs e as normais do objecto.


\begin{lstlisting}[caption=\textit{Buffers} funções.]
glGenBuffers(1, &a->vertexbuffer);
glBindBuffer(GL_ARRAY_BUFFER, a->vertexbuffer);
glBufferData(GL_ARRAY_BUFFER, (&Sphere->vertices)->size() * sizeof(glm::vec3), &Sphere->vertices[0][0], GL_STATIC_DRAW);

glGenBuffers(1, &a->uvbuffer);
glBindBuffer(GL_ARRAY_BUFFER, a->uvbuffer);
glBufferData(GL_ARRAY_BUFFER, (&Sphere->uvs)->size() * sizeof(glm::vec2), &Sphere->uvs[0][0], GL_STATIC_DRAW);
\end{lstlisting}

\subsubsection{\textit{draw(), setMVP()}}
\begin{lstlisting}[caption=\textit{setMVP()}.]
glm::mat4 setMVP(TRANSFORMATIONS b) {
	// Compute the MVP matrix from keyboard and mouse input
	computeMatricesFromInputs();
	glm::mat4 ProjectionMatrix = getProjectionMatrix();
	glm::mat4 ViewMatrix = getViewMatrix();
	glm::mat4 ModelMatrix = glm::mat4(1.0);
	glm::mat4 MVP = ProjectionMatrix * ViewMatrix * ModelMatrix * b.center_rotation * b.translation * b.self_rotation * b.scaling;
	return MVP;
}
\end{lstlisting}
Esta função vai receber como parâmetro a \textit{struct Transformations}, que vai ser usada para aplicar à matriz MVP (\textit{computeMatricesFromInputs()}). A matriz MVP foi calculada da seguinte forma : \textit{MVP = ProjectionMatrix * ViewMatrix * ModelMatrix * b.center\_rotation * b.translation * b.self\_rotation * b.scaling}.

\begin{lstlisting}[caption=\textit{draw()}.]
void draw(INFO *a, TRANSFORMATIONS b, GLuint programID, GLuint primitive, std::vector<glm::vec3>* verticesName) {
	// Use our shader
	
	glUseProgram(programID);
	glm::mat4 MVP= setMVP(b);
	// Send our transformation to the currently bound shader, 
	// in the "MVP" uniform
	glUniformMatrix4fv(a->MatrixID, 1, GL_FALSE, &MVP[0][0]);

	// Bind our texture in Texture Unit 0
	glActiveTexture(GL_TEXTURE0);
	glBindTexture(GL_TEXTURE_2D, a->Texture);
	// Set our "myTextureSampler" sampler to use Texture Unit 0
	glUniform1i(a->TextureID, 0);

	// 1rst attribute buffer : vertices
	glEnableVertexAttribArray(0);
	glBindBuffer(GL_ARRAY_BUFFER, a->vertexbuffer);
	glVertexAttribPointer(
		0,                  // attribute
		3,                  // size
		GL_FLOAT,           // type
		GL_FALSE,           // normalized?
		0,                  // stride
		(void*)0            // array buffer offset
	);

	// 2nd attribute buffer : UVs
	glEnableVertexAttribArray(1);
	glBindBuffer(GL_ARRAY_BUFFER, a->uvbuffer);
	glVertexAttribPointer(
		1,                                // attribute
		2,                                // size
		GL_FLOAT,                         // type
		GL_FALSE,                         // normalized?
		0,                                // stride
		(void*)0                          // array buffer offset
	);

	// Draw the triangle !
	glDrawArrays(primitive, 0, verticesName->size());

	glDisableVertexAttribArray(0);
	glDisableVertexAttribArray(1);
}
\end{lstlisting}

A função \textit{draw()} recebe como parâmetros as structs INFO e TRANSFORMATIONS, o \textit{programID} (link para os shaders), a \textit{primitive} e o \textit{verticesName} (vetor com os vertices).
A \textit{glActiveTexture} e a \textit{glBindTexture} são usadas para carregar as texturas.

A primeira ocorrência destas funções, \textit{glEnableVertexAttribArray}, \textit{glBindBuffer}, \textit{glVertexAttribPointer} envia os vértices deste objeto para os shaders.
A segunda, envia para os shaders os UVs, que por si são as coordenadas das texturas.

Posteriormente ele vai desenhar de acordo com a primitiva passada como parametro, desde o índice 0 até ao tamanho do vetor.


\subsubsection{\textit{resetPosPlanetsList()}}
Esta função vai dar \textit{reset} à posição da camâra quando o utilizador alterna entre planetas no Menu \textit{Planets}
\begin{lstlisting}[caption=\textit{resetPosPlanetsList()}.]
void resetPosPlanetsList() {
	position = glm::vec3(145.78f, 2.13f, 39.61f);
	horizontalAngle = -1.55f;
	verticalAngle= -0.03f;
\end{lstlisting}

\subsubsection{\textit{selectingMenu1()}}
Esta função é responsável pela escolha da opção do Menu Inicial.
\begin{lstlisting}[caption=\textit{ Pequeno excerto da função selectingMenu1()}.]
	if (glfwGetKey(window, GLFW_KEY_ENTER) == GLFW_PRESS && glfwGetKey(window, GLFW_KEY_ENTER) == GLFW_RELEASE)
	{
		if (selectedMenu1 == 1)
			currentMenu = 2;
		else if (selectedMenu1 == 2) {
			resetPosPlanetsList();
			currentMenu = 3;
		}

		else if (selectedMenu1 == 3)
			currentMenu = 4;
	}
\end{lstlisting}


\subsubsection{\textit{selectingMenu3()}}
Esta função é a responsável pela escolha da história dos planetas que irão aparecer.
\begin{lstlisting}[caption=\textit{ Pequeno excerto da função selectingMenu3()}.]
draw(&Background, transfBackground, programID, GL_TRIANGLES, &Back.vertices);
	if (glfwGetKey(window, GLFW_KEY_DOWN) == GLFW_PRESS && glfwGetKey(window, GLFW_KEY_DOWN) == GLFW_RELEASE)
	{
		resetPosPlanetsList();
		if (selectedMenu3 == 11)
			selectedMenu3 = 1;
		else
			selectedMenu3++;
	}
	if (glfwGetKey(window, GLFW_KEY_UP) == GLFW_PRESS && glfwGetKey(window, GLFW_KEY_UP) == GLFW_RELEASE)
	{
		resetPosPlanetsList();
		if (selectedMenu3 == 1)
			selectedMenu3 = 11;
		else
			selectedMenu3--;
	}
\end{lstlisting}

\subsubsection{\textit{declareDimensions()}}
Nesta função nós declaramos a escala real dos planetas e uma escala \textit{Fake}, sendo esta última usada para visualizar melhor os planetas. Atribuimos um valor \textit{random} ã posição inicial dos planetas em torno do Sol.
\newpage
\begin{lstlisting}[caption=\textit{Pequeno excerto da função declareDimensions()}.]
Real.sun = 0.0f * AU;
Real.mercury = 0.28f * AU;
Real.venus = 0.367f * AU;

(...)

Fake.sun = 0.0f * AU;
Fake.mercury = 0.28f * AU;
Fake.venus = 0.367f * AU;

(...)

UsingRotation.sun.center_Rotation = (float)(rand() % 365);
UsingRotation.mercury.center_Rotation = (float)(rand() % 365);
UsingRotation.venus.center_Rotation = (float)(rand() % 365);
\end{lstlisting}

\subsubsection{\textit{incrementRotations()}}
Esta função é a responsável pela escala de rotação e translação dso planetas.
\begin{lstlisting}[caption=\textit{ Pequeno excerto da função incrementRotations()}.]
UsingRotation.sun.center_Rotation += 0.0f;
UsingRotation.mercury.center_Rotation += 4.20f * rotationScale;
UsingRotation.venus.center_Rotation += 1.43f * rotationScale;

(...)

UsingRotation.sun.self_Rotation += 0.0115f;
UsingRotation.mercury.self_Rotation += 0.017f;
UsingRotation.venus.self_Rotation -= 0.004f;
\end{lstlisting}

\subsubsection{\textit{CelestialBodiesPosition()}}
Esta função define a translação em torno do Sol, a rotação sobre si mesma e o tamanho do planeta(\textit{scaling}).

\begin{lstlisting}[caption=\textit{ Pequeno excerto da função CelestialBodiesPosition()}.]
//Sol
	transfSun.center_rotation = glm::rotate(glm::mat4(1.0), glm::radians(UsingRotation.sun.center_Rotation), glm::vec3(0, 1, 0));
	transfSun.self_rotation = glm::rotate(glm::mat4(1.0), glm::radians(UsingRotation.sun.self_Rotation), glm::vec3(0, 1, 0));
	transfSun.translation = glm::translate(glm::mat4(1.0), glm::vec3(Using.sun, 0.0f, 0.0f));
	transfSun.scaling = glm::scale(glm::mat4(1.0), glm::vec3(19.0f));

//Mercury
	transfMercury.center_rotation = glm::rotate(glm::mat4(1.0), glm::radians(UsingRotation.mercury.center_Rotation), glm::vec3(0, 1, 0));
	transfMercury.self_rotation = glm::rotate(glm::mat4(1.0), glm::radians(UsingRotation.mercury.self_Rotation), glm::vec3(0, 1, 0));
	transfMercury.translation = glm::translate(glm::mat4(1.0), glm::vec3(Using.mercury, 0.0f, 0.0f));
\end{lstlisting}

\subsubsection{\textit{deletebuffers()}}
Nesta função apagamos o \textit{programID}, todas as texturas, \textit{buffers}(contêm vértices dos objectos e UVs). Apagamos ainda os \textit{VertexArrayID}'s.
\begin{lstlisting}[caption=\textit{ Pequeno excerto da função deletebuffers()}.]
glDeleteProgram(programID);

//Delete Textures
glDeleteTextures(1, &Background.Texture);
glDeleteTextures(1, &Sun.Texture);

//Delete buffers
glDeleteBuffers(1, &Background.vertexbuffer);
glDeleteBuffers(1, &Background.uvbuffer);

//Delete VertexArrays
glDeleteVertexArrays(1, &Background.VertexArrayID);
glDeleteVertexArrays(1, &Sun.VertexArrayID);
\end{lstlisting}

\subsubsection{\textit{Cintura de Asteróides}}
Depois de declarmos o randomAsteroid[100]  do tipo \textit{strut} TRANSFORMATIONS, percorremos cada posição do Array e atribuímos uma posição para o asteróide correspondente e um tamanho variável. De forma a tornar isto o mais \textit{random} possível gerámos uma nova \textit{seed} em cada execução do programa.

\begin{lstlisting}[caption=\textit{ Excerto de código para a criação da cintura de Asteroides}.]
    int j = 0;
    float i = 0;
    //Random seed
    time_t t;
    srand((unsigned) time(&t));
    for (i = 0, j = 0; i < 360; i += 3.6, j++)
    {
        //float cent_Rotation = (((float)(rand()) / RAND_MAX) * i) + i;
        float cent_Rotation = (rand() % 4) + i - 0.4;
        float transl = ((rand() % 10) + AU - 10);
        float scal = ((float)(rand()) / RAND_MAX) * 0.1f;

        randomAsteroid[j].self_rotation = glm::rotate(glm::mat4(1.0), glm::radians(1.0f), glm::vec3(0, 1, 0));
        randomAsteroid[j].center_rotation = glm::rotate(glm::mat4(1.0), glm::radians(cent_Rotation), glm::vec3(0, 1, 0));
        randomAsteroid[j].translation = glm::translate(glm::mat4(1.0), glm::vec3(transl, 0.0f, 0.0f));
        randomAsteroid[j].scaling = glm::scale(glm::mat4(1.0), glm::vec3(scal));
    }
\end{lstlisting}

\begin{figure}[H]
    \centering
    \includegraphics[width=12cm]{Images/cinturaasteroides.png}
    \caption{Cintura de Asteróides vista de cima}
    \label{fig:label}
\end{figure}


\begin{figure}[H]
    \centering
    \includegraphics[width=12cm]{Images/Menu.png}
    \caption{Menu Inicial}
    \label{fig:label}
\end{figure}

\begin{figure}[H]
    \centering
    \includegraphics[width=12cm]{Images/SolarSystemPerspective.png}
    \caption{Perspetiva do Sistema Solar}
    \label{fig:label}
\end{figure}

\begin{figure}[H]
    \centering
    \includegraphics[width=12cm]{Images/earth_history.png}
    \caption{História da Terra}
    \label{fig:label}
\end{figure}

\section{\textit{Shaders}}
\subsection{\textit{VertexShader}}
Os \textit{shaders} recebem as posições dos vértices do objeto que são armazenadas no \textit{vertexPosition\_modelspace}. Recebem também as posições dos objetos no qual terá a textura, estas são guardadas no \textit{vertexUV}.

O MVP é uma matriz 4x4 recebida do \textit{draw()} através do: \newline \textit{glUniformMatrix4fv(a->MatrixID, 1, GL\_FALSE, \&MVP[0][0])}. 

Primeiramente transformamos o \textit{vertexPosition\_modelspace} num vec4, uma vez que a matriz MVP é um mat4. Desta forma a multiplicação é possível para obter o \textit{gl\_Position}.

O UV recebe os \textit{VertexUV} e envia-os para o \textit{FragmentShader}.

\begin{lstlisting}[caption=\textit{VertexShader}]
#version 330 core

layout(location = 0) in vec3 vertexPosition_modelspace;
layout(location = 1) in vec2 vertexUV;

out vec2 UV;

uniform mat4 MVP;

void main(){

	gl_Position =  MVP * vec4(vertexPosition_modelspace,1);
	
	UV = vertexUV;
}
\end{lstlisting}

\subsection{\textit{FragmentShader}}
O \textit{FragmentShader} recebe os UV's vindos do \textit{VertexShader} e a textura (\textit{myTextureSampler}) recebida do \textit{transferDataToGPU()} através da função:\newline
a->TextureID = glGetUniformLocation(programID, "myTextureSampler").

Aplicamos a função \textit{texture()} onde passamos como argumento a textura e os UV's. Esta tranformação é enviada como vec3 para o OpenGL.

Como só existe um único \textit{out} o OpenGL já sabe qual é a variável das cores.

\begin{lstlisting}[caption=\textit{FragmentShader}]
#version 330 core

in vec2 UV;

out vec3 color;

uniform sampler2D myTextureSampler;

void main(){

	color = texture( myTextureSampler, UV ).rgb;
}
\end{lstlisting}





\clearpage{\thispagestyle{empty}\cleardoublepage}
\chapter{Conclusões e Trabalho Futuro}
\label{chap:conc-trab-futuro}

\section{Conclusão}
\label{sec:conc-princ}
A concretização do projeto  foi um processo demorado, mas simultanemante estimulante e desafiador para os alunos. A  realização  deste  projeto  constituiu-se  uma  mais  valia  na nossa formação e espelha a sistematização dos conhecimentos teórico-práticos adquiridos ao longo da Unidade Curricular.

Permitiu-nos expandir e consolidar os nossos conhecimentos referentes a C++ e ao OpenGL.

Devido à falta de tempo, não conseguimos implementar o sistema de iluminação como desejávamos, sendo algo a ter em conta em versões futuras do programa.

Um ponto a destacar foi a oportunidade de usar o Blender que facilitou a construção dos objetos.
Foi conseguida a interação do sistema teclado e rato para com o utilizador. As texturas deram um realce mais estético ao nosso Sistema Solar. A implementação de uma funcionalidade no Menu que permite consultar a história de um planeta à escolha, permite ao utilizador aumentar os seus conhecimentos.  

Posto isto e tendo o prazo sido cumprido, achamos que o nosso projeto foi concluido com sucesso.


\section{Trabalho Futuro}
\label{sec:trab-futuro}

No futuro pretendemos aplicar uma melhor iluminação e seria nossa intenção explorar formas de otimizar melhor o código. Seria de salientar a implementação de novas funcionalidades ao programa bem como melhorar a interação com o \textit{user}. 

\begin{thebibliography}{9}
\bibitem{Learn OpenGL} 
\textit{https://learnopengl.com/}

\bibitem{Página da disciplina} 
\textit{https://www.di.ubi.pt/~agomes/cg/}

\bibitem{Solar System Scope} 
\textit{https://www.solarsystemscope.com/textures/}

\bibitem{NASA Solar System Exploration} 
\textit{https://solarsystem.nasa.gov/}

\bibitem{Scale Model of Solar System} 
\textit{https://www.education.com/science-fair/article/scale-model-planets-solar-system/}

\bibitem{Explratorium} 
\textit{https://www.exploratorium.edu/ronh/solar_system/scale.pdf}
\end{thebibliography}
\clearpage{\thispagestyle{empty}\cleardoublepage}

% SE EXISTIREM APENDICES, DESCOMENTAR O QUE ESTÁ EM BAIXO
% \appendix
% \include{apendice1}
% \clearpage{\pagestyle{empty}\cleardoublepage}
% \include{continuacao}
% \clearpage{\pagestyle{empty}\cleardoublepage}
% \include{apendice2}
% \clearpage{\pagestyle{empty}\cleardoublepage}
% \include{apendice3}
% \clearpage{\pagestyle{empty}\cleardoublepage}

\backmatter

\bibliographystyle{unsrt}
%\bibliography{bibliografia}

\end{document}